% --- Archon physics formalization source begin ---
% archon:physics
% archon:covers PhyXMiniProblems/problem_phyx_mini_0045.lean
% archon:source-report reports/phyx_mini/problem_phyx_mini_0045.source.json

\chapter{Physics problem phyx\_mini\_0045}
\label{ch:PhyXMiniProblems_problem_phyx_mini_0045}

\paragraph{Problem source.}
\#\# Physical scenario

The far point of a certain myopic eye is 50 cm in front of the eye. Assume that the lens is worn 2 cm in front of the eye.

\#\# Question

Find the focal length of the eyeglass lens that will permit the wearer to see clearly an object at infinity.

\#\# Answer choices

- A: A: -35cm

- B: B: -48cm

- C: C: -37cm

- D: D: -43cm

\#\# Figure caption (auxiliary; use the image as primary evidence)

The image is a diagram illustrating the behavior of light rays passing through a diverging lens. 

- **Objects**: 
  - A horizontal line represents the optical axis.
  - A diverging lens is positioned along the axis.
  - To the left of the lens, a line labeled "Object at infinity" is shown with a zigzag indicating the object is infinitely far away.
  - An eye is drawn in dashed lines to the right of the lens, suggesting observation of the rays after passing through the lens.

- **Rays**:
  - A set of parallel purple rays is shown entering the lens from the left.
  - Upon passing through the lens, the rays diverge outward.
  - A dashed line extends back from the lens to where the rays appear to originate from, converging at a point labeled with a dot.

- **Measurements**:
  - The distances are labeled as \textbackslash{}( s = \textbackslash{}infty \textbackslash{}) for the object distance and \textbackslash{}( s' = f = -48 \textbackslash{}, \textbackslash{}text\{cm\} \textbackslash{}) for the image distance, indicating that the image is formed at the focal point on the same side as the object.

- **Text**:
  - A description in blue states, "When the object distance is infinity, all rays are parallel to the axis and the image distance equals the focal distance."
  - The label "Diverging lens" is written near the lens. 

This diagram shows how parallel rays are spread out by a diverging lens, causing the focal point to appear as though it's on the side of the lens where the rays originated.

\#\# Dataset metadata

Geometrical Optics; Physical Model Grounding Reasoning, Multi-Formula Reasoning

\paragraph{Recorded answer/context.}
B: B: -48cm

\paragraph{Figure/image path.}
/root/proposal\_for\_physic/science-mango/phyx\_mini\_run/phyx\_data/test\_image/45.png

\paragraph{Formalization target.}
create a compiling Lean file with sorry bodies at `PhyXMiniProblems/problem\_phyx\_mini\_0045.lean`. The Lean declarations must preserve the physical quantities, dimensions or dimensional roles, figure labels, governing-law hypotheses, and final relation expressed by this problem.
Use Mathlib/Physlib names found through LeanExplore where available. If a physics API is missing, introduce faithful local abstractions rather than scalar placeholder aliases.

\begin{theorem}[Physics formalization target]
\label{thm:physics:phyx_mini_0045:target}
\lean{PhyXMiniProblems.ProblemPhyXMini0045.problem_phyx_mini_0045}
\uses{def:physics:phyx-mini-0045:phyxminiproblems-problemphyxmini0045-lengthincentimeters, def:physics:phyx-mini-0045:phyxminiproblems-problemphyxmini0045-eyeglasscorrectionsetup, def:physics:phyx-mini-0045:phyxminiproblems-problemphyxmini0045-hasstateddistancereadouts, def:physics:phyx-mini-0045:phyxminiproblems-problemphyxmini0045-matchesprimaryfigure, def:physics:phyx-mini-0045:phyxminiproblems-problemphyxmini0045-hasphysicalaxialplacement, def:physics:phyx-mini-0045:phyxminiproblems-problemphyxmini0045-formsvirtualimageatfarpoint, def:physics:phyx-mini-0045:phyxminiproblems-problemphyxmini0045-obeysparallelrayfocallaw, def:physics:phyx-mini-0045:phyxminiproblems-problemphyxmini0045-answerfocallengthincentimeters}
The assigned autoformalize agent should translate this physics problem into Lean declarations in the covered file, with theorem and lemma proof bodies written as `by sorry`.
\end{theorem}
\begin{proof}
This is an autoformalization task, not a proof task. Produce statements that can later be proved by the physics prover without weakening the physical model.
\end{proof}

% --- Archon declaration topology begin ---
\section{Declaration topology}
\begin{definition}[Length Quantity]
  \label{def:physics:phyx-mini-0045:phyxminiproblems-problemphyxmini0045-lengthquantity}
  \lean{PhyXMiniProblems.ProblemPhyXMini0045.LengthQuantity}
  A signed physical length, independent of the unit chosen to read it
\end{definition}
\begin{proof}
  This is the declared model definition of Length Quantity.
\end{proof}
\begin{definition}[centimeter Unit Choices]
  \label{def:physics:phyx-mini-0045:phyxminiproblems-problemphyxmini0045-centimeterunitchoices}
  \lean{PhyXMiniProblems.ProblemPhyXMini0045.centimeterUnitChoices}
  Unit choices whose length component is the centimeter used in the source
\end{definition}
\begin{proof}
  This is the declared model definition of centimeter Unit Choices.
\end{proof}
\begin{definition}[length In Centimeters]
  \label{def:physics:phyx-mini-0045:phyxminiproblems-problemphyxmini0045-lengthincentimeters}
  \lean{PhyXMiniProblems.ProblemPhyXMini0045.lengthInCentimeters}
  \uses{def:physics:phyx-mini-0045:phyxminiproblems-problemphyxmini0045-lengthquantity, def:physics:phyx-mini-0045:phyxminiproblems-problemphyxmini0045-centimeterunitchoices}
  The signed scalar centimeter readout of a dimensionful length
\end{definition}
\begin{proof}
  This is the declared model definition of length In Centimeters.
\end{proof}
\begin{definition}[Thin Lens Kind]
  \label{def:physics:phyx-mini-0045:phyxminiproblems-problemphyxmini0045-thinlenskind}
  \lean{PhyXMiniProblems.ProblemPhyXMini0045.ThinLensKind}
  The two thin-lens types distinguished by the sign of focal length
\end{definition}
\begin{proof}
  This is the declared model definition of Thin Lens Kind.
\end{proof}
\begin{definition}[Eyeglass Lens]
  \label{def:physics:phyx-mini-0045:phyxminiproblems-problemphyxmini0045-eyeglasslens}
  \lean{PhyXMiniProblems.ProblemPhyXMini0045.EyeglassLens}
  \uses{def:physics:phyx-mini-0045:phyxminiproblems-problemphyxmini0045-lengthquantity, def:physics:phyx-mini-0045:phyxminiproblems-problemphyxmini0045-thinlenskind}
  A paraxial eyeglass lens represented by its signed physical focal length
\end{definition}
\begin{proof}
  This is the declared model definition of Eyeglass Lens.
\end{proof}
\begin{definition}[Object Location]
  \label{def:physics:phyx-mini-0045:phyxminiproblems-problemphyxmini0045-objectlocation}
  \lean{PhyXMiniProblems.ProblemPhyXMini0045.ObjectLocation}
  \uses{def:physics:phyx-mini-0045:phyxminiproblems-problemphyxmini0045-lengthquantity}
  The distance-labeled optical states needed for this problem
\end{definition}
\begin{proof}
  This is the declared model definition of Object Location.
\end{proof}
\begin{definition}[Incident Ray Geometry]
  \label{def:physics:phyx-mini-0045:phyxminiproblems-problemphyxmini0045-incidentraygeometry}
  \lean{PhyXMiniProblems.ProblemPhyXMini0045.IncidentRayGeometry}
  The relation of an incident ray bundle to the figure's optical axis
\end{definition}
\begin{proof}
  This is the declared model definition of Incident Ray Geometry.
\end{proof}
\begin{definition}[Myopic Eye]
  \label{def:physics:phyx-mini-0045:phyxminiproblems-problemphyxmini0045-myopiceye}
  \lean{PhyXMiniProblems.ProblemPhyXMini0045.MyopicEye}
  \uses{def:physics:phyx-mini-0045:phyxminiproblems-problemphyxmini0045-lengthquantity}
  A myopic eye together with its finite unaided far-point distance
\end{definition}
\begin{proof}
  This is the declared model definition of Myopic Eye.
\end{proof}
\begin{definition}[Eyeglass Correction Setup]
  \label{def:physics:phyx-mini-0045:phyxminiproblems-problemphyxmini0045-eyeglasscorrectionsetup}
  \lean{PhyXMiniProblems.ProblemPhyXMini0045.EyeglassCorrectionSetup}
  \uses{def:physics:phyx-mini-0045:phyxminiproblems-problemphyxmini0045-lengthquantity, def:physics:phyx-mini-0045:phyxminiproblems-problemphyxmini0045-eyeglasslens, def:physics:phyx-mini-0045:phyxminiproblems-problemphyxmini0045-objectlocation, def:physics:phyx-mini-0045:phyxminiproblems-problemphyxmini0045-incidentraygeometry, def:physics:phyx-mini-0045:phyxminiproblems-problemphyxmini0045-myopiceye}
  The physical objects and labeled distances in the eyeglass correction setup. signedVirtualImageDistanceFromLens is negative when the image is virtual and lies in front of the lens, on the same side as the object. The lens is a positive distance in front of the eye
\end{definition}
\begin{proof}
  This is the declared model definition of Eyeglass Correction Setup.
\end{proof}
\begin{definition}[Has Stated Distance Readouts]
  \label{def:physics:phyx-mini-0045:phyxminiproblems-problemphyxmini0045-hasstateddistancereadouts}
  \lean{PhyXMiniProblems.ProblemPhyXMini0045.HasStatedDistanceReadouts}
  \uses{def:physics:phyx-mini-0045:phyxminiproblems-problemphyxmini0045-lengthincentimeters, def:physics:phyx-mini-0045:phyxminiproblems-problemphyxmini0045-eyeglasscorrectionsetup}
  The two numerical distances explicitly supplied by the problem
\end{definition}
\begin{proof}
  This is the declared model definition of Has Stated Distance Readouts.
\end{proof}
\begin{definition}[Matches Primary Figure]
  \label{def:physics:phyx-mini-0045:phyxminiproblems-problemphyxmini0045-matchesprimaryfigure}
  \lean{PhyXMiniProblems.ProblemPhyXMini0045.MatchesPrimaryFigure}
  \uses{def:physics:phyx-mini-0045:phyxminiproblems-problemphyxmini0045-eyeglasscorrectionsetup}
  Qualitative readouts from the primary figure: the lens is diverging, the object is at infinity, and the incident rays are parallel to the optical axis
\end{definition}
\begin{proof}
  This is the declared model definition of Matches Primary Figure.
\end{proof}
\begin{definition}[Has Physical Axial Placement]
  \label{def:physics:phyx-mini-0045:phyxminiproblems-problemphyxmini0045-hasphysicalaxialplacement}
  \lean{PhyXMiniProblems.ProblemPhyXMini0045.HasPhysicalAxialPlacement}
  \uses{def:physics:phyx-mini-0045:phyxminiproblems-problemphyxmini0045-lengthincentimeters, def:physics:phyx-mini-0045:phyxminiproblems-problemphyxmini0045-eyeglasscorrectionsetup}
  The axial placement described in the text: both stated distances are positive and the spectacle lens lies between the far point and the eye
\end{definition}
\begin{proof}
  This is the declared model definition of Has Physical Axial Placement.
\end{proof}
\begin{definition}[Forms Virtual Image At Far Point]
  \label{def:physics:phyx-mini-0045:phyxminiproblems-problemphyxmini0045-formsvirtualimageatfarpoint}
  \lean{PhyXMiniProblems.ProblemPhyXMini0045.FormsVirtualImageAtFarPoint}
  \uses{def:physics:phyx-mini-0045:phyxminiproblems-problemphyxmini0045-eyeglasscorrectionsetup}
  The correction requirement: to be clear to the unaided myopic eye, the lens's virtual image must lie at the eye's far point. With distances from the eye given as positive magnitudes, its signed lens-relative distance is -(farPointDistance - lensToEyeDistance). The relation is required in every choice of units, so it is not tied to the centimeter readout
\end{definition}
\begin{proof}
  This is the declared model definition of Forms Virtual Image At Far Point.
\end{proof}
\begin{definition}[Obeys Parallel Ray Focal Law]
  \label{def:physics:phyx-mini-0045:phyxminiproblems-problemphyxmini0045-obeysparallelrayfocallaw}
  \lean{PhyXMiniProblems.ProblemPhyXMini0045.ObeysParallelRayFocalLaw}
  \uses{def:physics:phyx-mini-0045:phyxminiproblems-problemphyxmini0045-eyeglasscorrectionsetup}
  The paraxial thin-lens law in the infinite-object branch shown by the figure: when incident rays are parallel to the optical axis, their signed image distance equals the signed focal length. This governing law contains no numerical answer
\end{definition}
\begin{proof}
  This is the declared model definition of Obeys Parallel Ray Focal Law.
\end{proof}
\begin{definition}[Answer Choice]
  \label{def:physics:phyx-mini-0045:phyxminiproblems-problemphyxmini0045-answerchoice}
  \lean{PhyXMiniProblems.ProblemPhyXMini0045.AnswerChoice}
  The four answer-choice labels printed with the problem
\end{definition}
\begin{proof}
  This is the declared model definition of Answer Choice.
\end{proof}
\begin{definition}[answer Focal Length In Centimeters]
  \label{def:physics:phyx-mini-0045:phyxminiproblems-problemphyxmini0045-answerfocallengthincentimeters}
  \lean{PhyXMiniProblems.ProblemPhyXMini0045.answerFocalLengthInCentimeters}
  \uses{def:physics:phyx-mini-0045:phyxminiproblems-problemphyxmini0045-answerchoice}
  The signed focal-length readout printed beside each choice, in centimeters
\end{definition}
\begin{proof}
  This is the declared model definition of answer Focal Length In Centimeters.
\end{proof}
\begin{lemma}[virtual image distance eq neg 48]
  \label{lem:physics:phyx-mini-0045:phyxminiproblems-problemphyxmini0045-virtual-image-distance-eq-neg-48}
  \lean{PhyXMiniProblems.ProblemPhyXMini0045.virtual_image_distance_eq_neg_48}
  \uses{def:physics:phyx-mini-0045:phyxminiproblems-problemphyxmini0045-lengthincentimeters, def:physics:phyx-mini-0045:phyxminiproblems-problemphyxmini0045-eyeglasscorrectionsetup, def:physics:phyx-mini-0045:phyxminiproblems-problemphyxmini0045-hasstateddistancereadouts, def:physics:phyx-mini-0045:phyxminiproblems-problemphyxmini0045-formsvirtualimageatfarpoint}
  The supplied distances and clear-vision requirement place the virtual image 48 cm in front of the eyeglass lens, hence at signed distance -48 cm
\end{lemma}
\begin{proof}
  The conclusion follows from the governing relations and figure readouts named in the statement of virtual image distance eq neg 48.
\end{proof}
% --- Archon declaration topology end ---
% --- Archon physics formalization source end ---
